% ============================================================
% sec/02trabajos_relacionados.tex
%
% SECCIÓN CENTRAL DE LA ETAPA 1.
% Tres ejes + tabla comparativa obligatoria + brechas en
% apartado propio de cierre.
%
% ESTADO: el Eje B está desarrollado. Los ejes A y C dependen
% de los bundles de Daniel y Sebastián.
% ============================================================
\section{Trabajos Relacionados}
\label{sec:relacionados}

\pc{Párrafo de metodología de revisión: bases consultadas, cadenas de
búsqueda, criterios de inclusión y exclusión, ventana temporal, y cuántos
trabajos se retuvieron de cuántos revisados. Lo aporta el protocolo de
búsqueda del Eje C.}

\razon{Este párrafo convierte la revisión en un procedimiento auditable en
lugar de una lista de lecturas. Hodge \textit{et al.} organizan su revisión
declarando los temas que la estructuran \cite{hodge_winprediction_2021}, y
Grinsztajn \textit{et al.} enuncian criterios explícitos de inclusión de
conjuntos de datos \cite{grinsztajn_trees_2022}; ambas prácticas son el
precedente de exigir transparencia en el procedimiento de revisión.}

% ------------------------------------------------------------
% EJE A
% ------------------------------------------------------------
\subsection{Eje A: el problema y su dominio}
\label{subsec:ejea}

\pc{Tres a cinco párrafos sintetizando al menos siete trabajos arbitrados
sobre predicción de desenlace en juegos competitivos y evaluación de
estados en juegos de cartas. Corresponde al bundle de Daniel. Debe
agruparse por familia de enfoque y no por artículo.}

A continuación se resumen los trabajos de este eje que forman parte de la
revisión desarrollada en el presente documento, sin perjuicio de la
síntesis completa que corresponde a este apartado.

Hodge \textit{et al.} abordan la predicción del ganador en partidas
profesionales de un videojuego de arena multijugador utilizando información
del estado disponible mientras la partida transcurre
\cite{hodge_winprediction_2021}. Los autores organizan las características
empleadas en la literatura según el momento en que están disponibles
---previas al juego, durante el juego y posteriores al juego--- y recogen
la observación de que las características disponibles durante el juego
concentran la información más útil para la predicción. Su revisión reúne
dieciséis trabajos previos y advierte que la variación en tamaño,
composición y versión del juego entre los conjuntos de datos empleados
obliga a tratar con cautela la comparación de las exactitudes reportadas.

\razon{Esta advertencia debe recogerse explícitamente porque es el
precedente de dominio de la declaración de comparabilidad que exige la
Sección~\ref{sec:baseline}.}

Xenopoulos \textit{et al.} estiman la probabilidad de victoria por ronda en
un videojuego de disparos tácticos y comparan tres poblaciones de jugadores
con distinto nivel de habilidad \cite{xenopoulos_winprob_2022}. Su
representación del estado combina información global de la ronda con
información agregada por bando, estructura análoga a la que se adopta en
este trabajo. Los autores reportan que los modelos se encuentran bien
calibrados dentro de su propio entorno, y que el modelo entrenado sobre la
población de menor nivel presenta mayor pérdida logarítmica y menor
calidad de calibración al evaluarse sobre las poblaciones profesionales.

Saravanan y Guzdial estudian la evolución del conjunto de estrategias
dominantes de una comunidad competitiva ante cambios de balance del juego
\cite{saravanan_metadiscovery_2024}. El trabajo se sitúa en el videojuego
competitivo de la franquicia Pok\'emon y no en el juego de cartas
coleccionables, distinción que debe mantenerse presente al interpretar sus
resultados. Los autores definen dicho conjunto de estrategias dominantes
como el conocimiento que excede las reglas del juego y que se manifiesta en
los personajes y estrategias predominantes dentro de la base de jugadores,
y observan que cambia de forma continua.

\razon{La aclaración sobre cuál juego de la franquicia estudia este trabajo
es obligatoria en su primera mención. Presentarlo sin la distinción
atribuiría al artículo un alcance que no tiene.}

\noindent\textbf{Enfoques predominantes.} \pc{Respuesta directa a partir de
la síntesis completa del eje.}

\noindent\textbf{Aspectos considerados resueltos.} \pc{Respuesta directa.}

\noindent\textbf{Aspectos abiertos.} \pc{Respuesta directa. Este párrafo
alimenta el apartado de brechas.}

% ------------------------------------------------------------
% EJE B
% ------------------------------------------------------------
\subsection{Eje B: las técnicas del \textit{track} seleccionado}
\label{subsec:ejeb}

El \textit{track} seleccionado corresponde al aprendizaje automático
clásico sobre datos tabulares. Este apartado establece el estado actual de
dichos métodos, los supuestos bajo los cuales operan, y las condiciones que
determinan la calidad de las probabilidades que producen.

\subsubsection{Estado actual de los métodos tabulares}

Borisov \textit{et al.} presentan una revisión sistemática de los enfoques
de aprendizaje profundo para datos tabulares heterogéneos, junto con una
comparación experimental frente a métodos clásicos bajo un protocolo
uniforme \cite{borisov_tabular_2024}. Los autores organizan la literatura
en tres familias: métodos de transformación de datos, arquitecturas
especializadas ---modelos híbridos y basados en mecanismos de atención--- y
modelos de regularización. Su evaluación reporta que los conjuntos de
árboles con potenciación del gradiente obtienen los mejores resultados en
la mayoría de los conjuntos de datos considerados, y concluye que la
conveniencia de emplear técnicas actuales de aprendizaje profundo sobre
datos tabulares admite en general una respuesta negativa, particularmente
para conjuntos heterogéneos de tamaño reducido.

Grinsztajn \textit{et al.} amplían la evidencia con una comparación sobre
cuarenta y cinco conjuntos de datos, contabilizando el costo de la
selección de hiperparámetros, y reportan que los modelos basados en árboles
resultan superiores para todo presupuesto de búsqueda considerado
\cite{grinsztajn_trees_2022}.

\subsubsection{Supuestos bajo los que operan estos métodos}

Grinsztajn \textit{et al.} no se limitan a constatar la diferencia de
desempeño, sino que identifican tres propiedades de los datos tabulares que
la explican \cite{grinsztajn_trees_2022}. En primer lugar, las funciones
objetivo asociadas a estos conjuntos no son suaves, mientras que las redes
neuronales presentan un sesgo hacia soluciones suaves; los modelos basados
en árboles, que aprenden funciones constantes por tramos, no exhiben dicho
sesgo. En segundo lugar, los conjuntos tabulares contienen una proporción
apreciable de características no informativas, ante las cuales los modelos
neuronales de tipo perceptrón multicapa resultan menos robustos. En tercer
lugar, las características de un conjunto tabular poseen significado
individual, propiedad que los procedimientos invariantes ante rotaciones no
pueden aprovechar.

Estas tres propiedades resultan aplicables al problema abordado en este
trabajo. \pc{Un párrafo argumentando cada una sobre el dominio concreto:
la irregularidad de la función objetivo asociada a los umbrales duros del
juego, la presencia previsible de características poco informativas en el
diccionario, y el significado individual de las variables del estado. La
redacción de este párrafo requiere el diccionario de datos definitivo, que
corresponde a la Sección~\ref{sec:analisistecnico}.}

\razon{El enunciado exige que la justificación técnica del \textit{track}
se sostenga con literatura y no con opinión. Citar el resultado de Borisov
sin el mecanismo de Grinsztajn dejaría la justificación en el nivel de
«los árboles funcionan mejor»; incluir el mecanismo permite argumentar por
qué funcionan mejor en este problema en particular.}

\noindent\textbf{Límites del alcance de esta evidencia.} Borisov
\textit{et al.} reportan que, en el mayor de los conjuntos evaluados, con
características predominantemente continuas, una arquitectura basada en
mecanismos de atención supera a los métodos clásicos, y sugieren que para
conjuntos tabulares de gran tamaño con dichas características las
arquitecturas neuronales podrían presentar ventaja
\cite{borisov_tabular_2024}. Esta observación resulta pertinente dado el
volumen de observaciones disponible en el presente trabajo. Sin embargo, la
condición relevante no es el número de filas sino el número de unidades
estadísticamente independientes, distinción que se desarrolla en la
Sección~\ref{sec:metodologia}, y las variables del estado de juego son
heterogéneas y no predominantemente continuas.

Por su parte, Grinsztajn \textit{et al.} excluyen de su comparación los
conjuntos con datos faltantes, aquellos con variables categóricas de
cardinalidad superior a veinte, los conjuntos cuyas observaciones no son
independientes, y los conjuntos deterministas procedentes de juegos, por
considerar que estos últimos difieren de los conjuntos tabulares habituales
y ameritan estudio separado \cite{grinsztajn_trees_2022}. El problema
abordado en este trabajo no es determinista: el desenlace de una partida no
se deduce del estado intermedio, sino que depende del orden restante del
mazo, de la información oculta del oponente y de las decisiones
posteriores de ambos jugadores. Es precisamente esta condición la que
motiva estimar una probabilidad en lugar de deducir una etiqueta. Respecto
del supuesto de independencia, este trabajo adopta el episodio como unidad
independiente, según se detalla en la Sección~\ref{sec:metodologia}.

\razon{Ambos límites deben declararse en el texto y no omitirse. Son las
dos objeciones más directas que puede recibir la justificación del
\textit{track}, y responderlas de forma anticipada es preferible a que
surjan durante la sustentación.}

\subsubsection{Dificultad intrínseca de la tarea}

Brill \textit{et al.} examinan la dificultad de estimar la probabilidad de
victoria mediante un estudio de simulación en el cual la probabilidad
verdadera es calculable de forma exacta \cite{brill_winprob_2026}. Los
autores reportan que la estructura de dependencia propia de los datos
observacionales, en la cual múltiples observaciones comparten el mismo
desenlace, incrementa el sesgo y la varianza de los estimadores y reduce el
tamaño de muestra efectivo respecto del nominal. Reportan asimismo que
resulta preferible conservar la totalidad de las observaciones por partida
antes que retener una sola, dado que el conjunto completo aporta
información sobre la estructura del espacio de covariables. Finalmente,
observan que los intervalos de confianza construidos mediante remuestreo
presentan cobertura inferior a la nominal, incluso cuando el remuestreo
respeta la estructura de agrupamiento.

Los autores declaran de forma explícita que sus hallazgos no son
exclusivos de la aplicación estudiada, sino que resultan aplicables a
cualquier conjunto de datos en el cual la variable de resultado
corresponda al desenlace de una unidad y las observaciones consistan en
las unidades que conducen a dicho desenlace
\cite{brill_winprob_2026}. El conjunto de datos empleado en este trabajo
satisface dicha descripción.

\subsubsection{Fuga de información por dependencia entre observaciones}

Figueiredo y Mendes estudian el efecto de la partición aleatoria sobre
conjuntos de imágenes extraídas de material audiovisual, en los cuales las
observaciones procedentes de una misma escena presentan alta correlación
\cite{figueiredo_leakage_2024}. Los autores comparan la partición
aleatoria contra una partición que asigna agrupaciones completas a cada
subconjunto, y reportan que la primera produce una inflación apreciable de
las métricas de desempeño respecto de la partición agrupada de referencia.
Su trabajo pertenece al dominio de la visión por computador, por lo que se
considera aquí un análogo metodológico y no un resultado sobre datos de
juegos; el mecanismo subyacente ---la presencia de observaciones casi
idénticas a ambos lados de la partición--- es independiente del dominio.

\razon{La declaración de la brecha de dominio es obligatoria. Presentar un
resultado de visión por computador como si procediera del dominio de
juegos constituiría una atribución incorrecta.}

Cabe destacar una diferencia favorable al presente trabajo. La
contribución central de Figueiredo y Mendes consiste en un procedimiento
para inferir las agrupaciones, dado que las imágenes no vienen etiquetadas
por escena \cite{figueiredo_leakage_2024}. En el conjunto de datos
empleado aquí, la agrupación viene dada de forma explícita por el
identificador de episodio, lo que permite aplicar la partición agrupada de
manera exacta y verificable.

\subsubsection{Evaluación de la calidad de las probabilidades}

Silva Filho \textit{et al.} presentan una revisión de los métodos para
evaluar y mejorar las probabilidades predichas por un clasificador
\cite{silvafilho_calibration_2023}. Un modelo se considera calibrado
cuando las probabilidades que predice se corresponden con las frecuencias
observadas. Los autores exponen que las reglas de puntuación propias se
descomponen en un componente de calibración y uno de refinamiento,
descomposición que fundamenta el conjunto de métricas adoptado en este
trabajo: las métricas basadas en ordenamiento informan sobre el
refinamiento, las medidas de error de calibración sobre el componente
homónimo, y las reglas de puntuación propias agregan ambos. Los autores
señalan asimismo que los estimadores del error de calibración basados en
discretización dependen del número y la construcción de los intervalos
empleados.

Niculescu-Mizil y Caruana examinan las probabilidades producidas por
distintas familias de algoritmos y reportan que los métodos de máximo
margen, entre ellos los conjuntos de árboles con potenciación, desplazan la
masa de probabilidad alejándola de los extremos, produciendo una distorsión
de forma sigmoidea \cite{niculescumizil_probabilities_2005}. Reportan
igualmente que la regresión logística y los conjuntos por agregación
resultan bien calibrados sin corrección posterior, y que aplicar
recalibración sobre modelos ya calibrados puede deteriorar su desempeño.
Su análisis de curvas de aprendizaje establece que el método de
recalibración apropiado depende del tamaño del conjunto disponible para
dicho ajuste.

\razon{Esta referencia debe aparecer porque convierte la calibración de una
preocupación general en una predicción concreta sobre los modelos que este
trabajo va a emplear. Sin ella, la sección de métricas carecería de
justificación específica.}

Ojeda \textit{et al.} comparan de forma sistemática los enfoques de
recalibración disponibles para la estimación de probabilidades de dos
clases, evaluando entre otros aspectos su capacidad de generalizar
estimaciones precisas a poblaciones externas
\cite{ojeda_calibrating_2023}. Dimitriadis \textit{et al.} abordan la
inestabilidad de los diagramas de fiabilidad construidos mediante
discretización, y proponen un procedimiento que genera diagramas
reproducibles junto con una medida numérica de descalibración
\cite{dimitriadis_reliability_2021}.

\decidir{Las recomendaciones de Ojeda \textit{et al.} y las de
Niculescu-Mizil y Caruana no coinciden de forma inmediata respecto del
método de recalibración preferible. La resolución de esta diferencia
requiere revisar si el escenario evaluado por Ojeda \textit{et al.}
---transferencia a poblaciones externas--- corresponde al protocolo
temporal de este trabajo, en cuyo caso su recomendación tendría
precedencia. La decisión debe tomarse antes de la Etapa 2.}

% ------------------------------------------------------------
% EJE C
% ------------------------------------------------------------
\subsection{Eje C: el conjunto de datos}
\label{subsec:ejec}

\pc{Tres a cinco párrafos correspondientes al bundle de Sebastián. El
conjunto de datos empleado carece de literatura arbitrada, por lo que este
apartado debe aplicar la sustitución prevista: revisar trabajos con
conjuntos análogos del mismo dominio y justificar en qué medida son
comparables.}

\razon{La ausencia de literatura sobre el conjunto de datos debe declararse
de forma explícita en este apartado y recogerse después como brecha, en
lugar de omitirse. Declararla convierte una carencia en un hallazgo de la
revisión.}

\noindent\textbf{Métricas y protocolos reportados.} \pc{Qué mide la
literatura sobre conjuntos análogos y bajo qué partición.}

\noindent\textbf{Problemas conocidos.} \pc{Sesgos, fugas, limitaciones de
tamaño o de representatividad. Este apartado justifica varias decisiones de
la Sección~\ref{sec:analisistecnico}.}

% ------------------------------------------------------------
% TABLA COMPARATIVA
% ------------------------------------------------------------
\subsection{Análisis comparativo}
\label{subsec:comparativo}

\safeinput{tab/tabla_comparativa}

\pc{Dos a tres párrafos interpretando la tabla: tendencias observables,
comparaciones que no resultan legítimas y por qué, y limitaciones que se
repiten entre trabajos.}

\razon{El texto debe leer la tabla y no repetirla. Si el párrafo se limita
a enumerar las filas, la tabla no cumple su función. Hodge \textit{et al.}
ofrecen el precedente de publicar la comparación acompañada de la
advertencia sobre su comparabilidad \cite{hodge_winprediction_2021}.}

% ------------------------------------------------------------
% BRECHAS
% ------------------------------------------------------------
\subsection{Brechas identificadas}
\label{subsec:brechas}

\noindent\textbf{Carencias encontradas en la literatura.}

\begin{enumerate}
  \item \textbf{Ausencia de estimación de probabilidad de victoria en
        juegos de cartas coleccionables a partir del estado de juego.}
        \pc{Verificar y precisar el alcance de esta afirmación una vez
        completados los ejes A y C.}
  \item \textbf{Evaluación desigual de la calidad de las probabilidades en
        el dominio de juegos.} La literatura de predicción de desenlace
        revisada reporta predominantemente métricas de exactitud o de
        ordenamiento \cite{hodge_winprediction_2021}, con la excepción de
        Xenopoulos \textit{et al.}, quienes incorporan una medida de error
        de calibración \cite{xenopoulos_winprob_2022}.
  \item \textbf{Evaluación de métricas probabilísticas pendiente en la
        comparación de modelos tabulares.} Grinsztajn \textit{et al.}
        identifican explícitamente la evaluación de las predicciones
        probabilísticas como línea de trabajo no cubierta por su
        comparación \cite{grinsztajn_trees_2022}.
  \item \textbf{Ausencia de caracterización de la degradación de la
        calibración bajo desplazamiento temporal en el dominio de juegos.}
        Xenopoulos \textit{et al.} evalúan la transferencia entre
        poblaciones de distinto nivel de habilidad
        \cite{xenopoulos_winprob_2022}, escenario distinto del
        desplazamiento temporal de la población.
  \item \pc{Brecha correspondiente al conjunto de datos, aportada por el
        Eje C.}
  \item \pc{Brecha correspondiente al dominio, aportada por el Eje A.}
\end{enumerate}

\razon{La brecha número dos no puede formularse como ausencia absoluta. Al
menos un trabajo de la propia bibliografía reporta calibración, de modo que
afirmar que ningún trabajo lo hace sería contradicho por una referencia
citada en este mismo documento.}

\noindent\textbf{Brechas que este trabajo aborda.} El presente trabajo
aborda las brechas primera, segunda y cuarta. La primera, mediante la
construcción y evaluación de un estimador sobre un dominio no cubierto por
la literatura revisada. La segunda y la cuarta, mediante el reporte
separado de discriminación y calibración bajo tres protocolos de
generalización, uno de los cuales corresponde a desplazamiento temporal.

\noindent\textbf{Brechas que quedan fuera del alcance.} La tercera brecha
---la evaluación sistemática de métricas probabilísticas en la comparación
de modelos tabulares--- excede el alcance de este trabajo, que no constituye
un estudio comparativo de familias de modelos sobre múltiples conjuntos de
datos. \pc{Añadir las brechas fuera de alcance correspondientes a los ejes
A y C.}

\razon{Reconocer explícitamente las brechas que el trabajo no aborda es un
requisito del enunciado y no una debilidad. Un apartado de brechas que
declara abordarlas todas resulta menos creíble que uno que delimita su
alcance.}